\section{Model}

Time is continuous. Population grows at rate \(n\), and people work either in
final production or AI research:
\begin{equation}
 \dot N=nN,\qquad N=L+H.
 \label{eq:axm-labor-market}
\end{equation}
Here \(L\) always denotes production labor and \(H\) human research labor.

\subsection{Household and final production}

The representative household maximizes
\[
 \int_0^\infty e^{-\rho t}N_t\log(C_t/N_t)\,\dd t
\]
and owns physical capital and the AI developer. Its interior Euler equation in
aggregate variables is
\begin{equation}
 \frac{\dot C}{C}=n+\alpha\frac{Y}{K}-\delta-\rho.
 \label{eq:axm-euler}
\end{equation}
Competitive final-good firms produce
\begin{align}
 Y&=K^\alpha Z^{1-\alpha},\\
 Z&=\left[\omega_L L^{\frac{\sigma_{XL}-1}{\sigma_{XL}}}
 +\omega_X X^{\frac{\sigma_{XL}-1}{\sigma_{XL}}}\right]^{
 \frac{\sigma_{XL}}{\sigma_{XL}-1}},
 \qquad \omega_L+\omega_X=1.
 \label{eq:axm-final-production}
\end{align}
The Cobb--Douglas limit is \(Z=L^{\omega_L}X^{\omega_X}\). Define the CES
contribution share of AI services as
\begin{equation}
 s_X=\frac{\omega_X X^{(\sigma_{XL}-1)/\sigma_{XL}}}
 {\omega_L L^{(\sigma_{XL}-1)/\sigma_{XL}}
 +\omega_X X^{(\sigma_{XL}-1)/\sigma_{XL}}}.
 \label{eq:axm-sx}
\end{equation}
Factor prices then satisfy
\begin{equation}
 w=(1-\alpha)(1-s_X)\frac{Y}{L},\qquad
 p_X=(1-\alpha)s_X\frac{Y}{X}.
 \label{eq:axm-factor-prices}
\end{equation}

\Needspace{11\baselineskip}
\begin{remark}[Tasks, jobs, and labor services]
The inputs \(L\) and \(H\) are quantities of homogeneous labor services, not
numbers of tasks, occupations, or jobs. A job can combine several tasks, and AI
can replace some of them without eliminating the job. The model abstracts from
that margin. In addition, \(N=L+H\) imposes full employment: people can move
between final production and AI research, but the model contains neither
unemployment nor job creation and destruction. Consequently, \(\sigma_{XL}\)
measures substitution between aggregate labor and AI services in final
production. It is not a task-level elasticity or a direct measure of the number
of jobs displaced by AI. The model's labor-market predictions concern wages,
factor shares, and the allocation of labor between the two sectors.
\end{remark}

\subsection{Two uses of compute}

\(A\) indexes current AI capability: what one unit of compute can accomplish.
\(U\) is inference compute and \(M\) is research compute. Their service flows are
\begin{equation}
 X=AU,\qquad R=AM.
 \label{eq:axm-two-compute-services}
\end{equation}
\(R\) is only shorthand for automated research services. The paid inputs are
\(U\) and \(M\). Because both uses have the same constant unit cost, we choose
units of raw compute so that one unit costs one unit of final output. This is a
normalization, not an assumption that compute is economically free. To see this,
start from a common cost \(\xi>0\), define
\(\widetilde U=\xi U\), \(\widetilde M=\xi M\),
\(\widetilde A=A/\xi\), \(\widetilde\chi=\chi/\xi\), and, for the shadow
value used below, \(\widetilde q=\xi q\), and then drop tildes.
Both service flows and all resource costs are unchanged. A unit of either AI
service therefore has marginal cost \(1/A\). This separates the capability stock
\(A\), two raw compute inputs \(U,M\), and the two service flows \(X,AM\).

\subsection{AI research}

Human researchers and automated research services produce effective research:
\begin{align}
 E&=\left[\omega_H H^{\frac{\sigma_{HM}-1}{\sigma_{HM}}}
 +\omega_M(AM)^{\frac{\sigma_{HM}-1}{\sigma_{HM}}}\right]^{
 \frac{\sigma_{HM}}{\sigma_{HM}-1}},
 \qquad \omega_H+\omega_M=1,\\
 \dot A&=\chi E^\eta,\qquad 0<\eta<1.
 \label{eq:axm-capability-law}
\end{align}
The baseline contains no additional \(A^\phi\) term. Better capability affects
research only by raising the services produced by \(M\). This isolates research
automation and avoids counting the capability feedback twice.

The unit cost of effective research is
\begin{equation}
 P_E=\left[\omega_H^{\sigma_{HM}}w^{1-\sigma_{HM}}
 +\omega_M^{\sigma_{HM}}\left(\frac{1}{A}\right)^{1-\sigma_{HM}}
 \right]^{1/(1-\sigma_{HM})}.
 \label{eq:axm-research-price}
\end{equation}
Conditional input demands are
\begin{equation}
 H=\omega_H^{\sigma_{HM}}\left(\frac{P_E}{w}\right)^{\sigma_{HM}}E,
 \quad
 AM=\omega_M^{\sigma_{HM}}
 \left(AP_E\right)^{\sigma_{HM}}E.
 \label{eq:axm-research-demands}
\end{equation}
The automated contribution share is
\begin{equation}
 s_M=\frac{M}{wH+M}.
 \label{eq:axm-sm}
\end{equation}
By constant returns in \(E\), it is also the elasticity of \(E\) with respect
to \(AM\).

\subsection{Integrated AI developer}

The developer supplies \(X\), pays \(X/A\) for inference compute, hires
\(H\), uses \(M\), and owns \(A\). The inverse elasticity of demand for \(X\) is
\begin{equation}
 \varepsilon_X^{-1}=\frac{1-s_X}{\sigma_{XL}}+\alpha s_X.
 \label{eq:axm-inverse-elasticity}
\end{equation}
On an interior monopoly branch,
\begin{equation}
 p_X(1-\varepsilon_X^{-1})=\frac{1}{A}.
 \label{eq:axm-monopoly-foc}
\end{equation}
Let \(q\) be the current-value shadow price of capability. Research optimality is
\begin{equation}
 q\frac{\partial\dot A}{\partial H}=w,\qquad
 q\frac{\partial\dot A}{\partial M}=1,
 \label{eq:axm-research-focs}
\end{equation}
or, after minimizing the cost of \(E\),
\begin{equation}
 q\chi\eta E^{\eta-1}=P_E.
 \label{eq:axm-research-scale-foc}
\end{equation}
The envelope derivative of operating profit is \(\pi_A=X/A^2\). Moreover,
\(\partial\dot A/\partial A=\eta s_M\dot A/A\) because a fixed \(M\) supplies
more research services when \(A\) rises. The costate equation is
\begin{equation}
 \frac{\dot q}{q}=\alpha\frac{Y}{K}-\delta-\frac{X}{qA^2}
 -\eta s_M\frac{\dot A}{A}.
 \label{eq:axm-costate}
\end{equation}
The last term is the main dynamic difference from research that combines \(H\)
directly with \(M\).

\subsection{Resources and equilibrium}

The resource constraint is
\begin{equation}
 C+\dot K+\delta K+U+M=Y.
 \label{eq:axm-resource}
\end{equation}

\begin{definition}[Equilibrium]
Given \(K(0),A(0),N(0)>0\), an equilibrium consists of paths satisfying household
optimality, competitive final-good demands, the monopoly and research first-order
conditions, the developer's costate and transversality conditions, both
technologies, and the labor and goods market-clearing conditions.
\end{definition}

The system used for computation is explicit. Given \(K,A,N,q\), equations
\eqref{eq:axm-labor-market} and
\eqref{eq:axm-final-production}--\eqref{eq:axm-research-scale-foc} determine the
static allocation. The dynamics are
\begin{align}
 \frac{\dot K}{K}&=\frac{Y-C-U-M}{K}-\delta,\\
 \frac{\dot A}{A}&=\frac{\chi E^\eta}{A},\\
 \frac{\dot C}{C}&=n+\alpha\frac{Y}{K}-\delta-\rho,\\
 \frac{\dot q}{q}&=\alpha\frac{Y}{K}-\delta-\frac{X}{qA^2}
 -\eta s_M\frac{\dot A}{A}.
 \label{eq:axm-canonical-system}
\end{align}
\(K\) and \(A\) are predetermined; \(C\) and \(q\) jump to the equilibrium saddle
path. The block is closed by initial \(K,A,N\), the labor and goods market
identities, and the no-Ponzi and transversality conditions. In current-value
notation, the relevant terminal restrictions include
\begin{equation}
 \lim_{t\to\infty}e^{-\rho t}\frac{N_tK_t}{C_t}=0,\qquad
 \lim_{t\to\infty}
 \exp\left\{-\int_0^t\left(\alpha\frac{Y_s}{K_s}-\delta\right)\dd s\right\}
 q_tA_t=0.
 \label{eq:axm-tvc}
\end{equation}
The numerical boundary conditions select the stable path approaching the
analytically characterized limit and are checked against these restrictions.

\begin{definition}[Balanced growth and singularity]
A balanced-growth path has finite constant growth rates and constant limiting
allocation shares. An accelerating path has an output growth rate that does not
converge to a finite constant. Following \citet{aghionjonesjones2019}, a
finite-time singularity occurs when output becomes unbounded at a finite date.
\end{definition}
